\chapter{Software-Defined Storage and Distributed Block Storage}
\label{ch:7}

The evolution from monolithic storage arrays to distributed, software-defined storage systems represents one of the most significant architectural shifts in enterprise infrastructure over the past two decades. Where traditional storage arrays coupled proprietary hardware with embedded firmware into vertically integrated appliances, distributed storage unbundles those functions: placing commodity servers in service as storage nodes, managing data across them through sophisticated software layers, and scaling by adding nodes rather than by replacing hardware. This chapter examines the foundational theory that governs distributed storage behavior, explores the major parallel filesystem and software-defined storage platforms, analyzes the scale-out NAS landscape, and dissects the data placement and metadata management algorithms that make large-scale distributed storage systems practical.

\section{The CAP Theorem and Distributed Storage Fundamentals}
\section{Understanding CAP in the Storage Context}
Eric Brewer's CAP theorem, formally proven by Gilbert and Lynch in 2002, states that a distributed data system cannot simultaneously guarantee all three of the following properties: Consistency (every read receives the most recent write or an error), Availability (every request receives a response, even if it may not be the most recent data), and Partition Tolerance (the system continues operating even when network partitions prevent some nodes from communicating with others).

For practitioners designing or selecting storage systems, the theorem's most important implication is not the impossibility of achieving all three simultaneously --- network partitions are an empirical reality in any distributed system, so partition tolerance is effectively non-negotiable --- but rather the forced trade-off between consistency and availability during a partition event. A storage system that prioritizes consistency during a partition will refuse to serve reads or writes until the partition heals and data can be verified coherent across all replicas. A system that prioritizes availability will continue serving requests but may return stale data or permit diverging writes that must later be reconciled.

In practice, distributed storage systems occupy a spectrum rather than a binary position. Most enterprise distributed storage systems are classified as CP (consistency over availability) or AP (availability over consistency), but the reality is more nuanced. Ceph, for example, makes pools unavailable when a quorum of OSDs cannot be reached, prioritizing consistency. GlusterFS can be configured for either consistency or eventual consistency depending on the translator stack. Many scale-out NAS systems maintain strong consistency within a cluster but offer eventual consistency for geographically distributed global namespaces.

\section{Consistency Models in Distributed Storage}
Beyond the binary of CAP, distributed storage systems offer a range of consistency guarantees that practitioners must understand when matching a storage platform to its intended workload.

\textbf{Strong consistency} (also called linearizability) guarantees that all reads reflect the most recent write, as if the storage were a single atomic unit. Systems providing strong consistency --- typically CP systems --- accomplish this through distributed locking, two-phase commit protocols, or quorum-based reads and writes. The performance cost is real: every write operation may require coordination across multiple nodes before being acknowledged, and reads may require a consistency check before being served.

\textbf{Sequential consistency} is a slightly weaker model where all operations appear to have executed in some sequential order, consistent across all nodes, but that order need not match real-time ordering. This model is sufficient for many applications and less expensive to implement than linearizability.

\textbf{Causal consistency} preserves ordering among related operations (operations that have a cause-and-effect relationship), while permitting unrelated operations to be observed in different orders by different nodes. Some distributed databases and eventually some distributed filesystems have adopted causal consistency as a practical middle ground.

\textbf{Eventual consistency} guarantees that, in the absence of further writes, all replicas will eventually converge to the same value. The "eventually" is undefined and can range from milliseconds in a well-connected LAN to minutes or longer in a geographically distributed system. Eventual consistency dramatically simplifies the write path --- a write is acknowledged as soon as it lands on one or more replicas, and background processes propagate changes. The exposure to stale reads is the trade-off.

\textbf{Read-your-writes consistency} is a client-side session guarantee: a client that writes a value will always read back that same value in subsequent reads, even if other clients may not yet see it. Many distributed file protocols implement this guarantee because applications almost universally expect it.

\section{PACELC: Extending the Model}
The PACELC theorem, proposed by Daniel Abadi, extends CAP by recognizing that the consistency-availability trade-off is not limited to partition scenarios. Even in the absence of a partition --- the normal operating condition --- there is an inherent trade-off between latency and consistency. A system that synchronously replicates a write to three nodes before acknowledging it is more consistent but slower than one that acknowledges after writing to one node and replicates asynchronously. PACELC formalizes this: Partition $\rightarrow$ tradeoff between Availability and Consistency; Else $\rightarrow$ tradeoff between Latency and Consistency.

For enterprise architects, PACELC provides a more complete framework for evaluating storage systems. A high-frequency trading workload may demand EL (low latency) even at the cost of potential inconsistency windows. A financial ledger demands EC (strong consistency) even if it means higher latency on each write. Most enterprise storage systems position themselves on this spectrum implicitly, and understanding where each product sits helps avoid painful mismatches between application requirements and storage behavior.

\section{Parallel Filesystems for High-Performance Computing}
Parallel filesystems were developed specifically to address the I/O demands of high-performance computing (HPC) workloads --- large-scale scientific simulations, genomics pipelines, seismic analysis, and weather modeling --- where aggregate bandwidth measured in hundreds of gigabytes per second is required and sequential large-file I/O dominates. These systems take a fundamentally different approach from general-purpose NAS, separating metadata management from data movement and enabling simultaneous access from thousands of compute nodes to shared files.

\section{IBM Spectrum Scale (GPFS)}
IBM General Parallel File System, now marketed as IBM Spectrum Scale, was introduced in the early 1990s and has evolved into one of the most capable and widely deployed parallel filesystems in the HPC and enterprise space. Its architecture is notable for using a distributed metadata approach: multiple metadata servers (called manager nodes) share responsibility for namespace operations, avoiding the single metadata server bottleneck that afflicts simpler parallel filesystems.

Spectrum Scale organizes storage as a collection of Network Shared Disks (NSDs) --- logical storage devices that can be backed by local disks, SAN volumes, or other storage subsystems --- which are pooled into storage pools and assigned to filesets (subtree namespaces that can be independently managed, migrated, and governed). The filesystem distributes data across NSDs in stripes, with configurable stripe size and replication factor per fileset.

The \textbf{distributed token management} mechanism is central to Spectrum Scale's consistency model. Tokens are held by nodes for specific operations (read, write, metadata) and must be obtained from a token manager before proceeding. This prevents conflicting operations from proceeding simultaneously across different nodes, ensuring coherent shared access. Token management adds latency for highly concurrent small-file workloads but is essentially transparent for the large, sequential I/O patterns that Spectrum Scale is optimized for.

Spectrum Scale's \textbf{Active File Management (AFM)} feature implements a global namespace with cache semantics: remote data can appear locally mounted, with background prefetch and writeback. This is particularly valuable for hybrid cloud scenarios where data may reside in cloud object storage (via cloud gateway) and needs to appear as a local filesystem to compute jobs.

The product's strength lies in its flexibility and enterprise feature set: per-fileset encryption, immutable file support, Hadoop HDFS compatibility mode, transparent tiering to object storage, and integration with IBM Spectrum Protect for data management. Spectrum Scale is typically found in financial services, oil and gas, healthcare genomics, and research computing environments requiring petabyte-scale shared storage with enterprise governance.

\section{Lustre}
Lustre is an open-source parallel filesystem originally developed by Cluster File Systems and later contributed to the open-source community. It is the dominant parallel filesystem in TOP500 supercomputing installations and is the reference platform for extreme-bandwidth HPC workloads. Its architecture separates three distinct server roles.

\textbf{Metadata Servers (MDS)} with their \textbf{Metadata Targets (MDT)} handle all namespace operations: directory lookups, file creates, stat operations, file open/close. In traditional Lustre deployments, a single active MDS was a significant architectural limitation, representing both a single point of failure and a performance ceiling for metadata-intensive workloads. Lustre 2.4 introduced \textbf{Distributed Namespace (DNE)}, allowing multiple MDTs and distributing directory subtrees across them, with subsequent releases adding striped directories that can shard even a single directory across multiple MDTs for workloads generating millions of small files.

\textbf{Object Storage Servers (OSS)} with their \textbf{Object Storage Targets (OST)} hold the actual file data. Files are striped across multiple OSTs, with stripe count and stripe size configured per file or directory. A file's data stripe is interleaved across its OSTs at the configured stripe width. With hundreds of OSTs and a well-tuned stripe configuration, Lustre deployments at major supercomputing centers regularly achieve aggregate bandwidths exceeding 1 TB/s.

\textbf{Management Servers (MGS)} maintain configuration information for one or more Lustre filesystems and are typically co-located with the MDS.

The \textbf{Lustre Distributed Lock Manager (LDLM)} provides cache coherence across all clients. When a client reads a file, it acquires a lock on the data extent from the MDS. Other clients attempting to modify that extent must first revoke the existing lock. This mechanism enables aggressive client-side read caching while preserving coherence.

Lustre's weakness has historically been in metadata performance for workloads generating large numbers of small files --- the classic HPC checkpoint problem --- and in operational complexity. Managing Lustre at scale requires expertise that many enterprise IT teams lack, which has contributed to the rise of more operationally friendly alternatives.

\section{BeeGFS}
BeeGFS (previously FhGFS, developed by the Fraunhofer Institute for Industrial Mathematics) takes a pragmatic approach to parallel filesystem design, optimizing for operational simplicity and excellent performance across mixed workload types while remaining deployable on commodity hardware without specialized expertise.

BeeGFS fully distributes both data and metadata across all storage nodes, with no architectural separation between metadata and data services --- any node can run both a metadata daemon and a storage daemon. Metadata is distributed across all metadata nodes using a consistent hashing approach, assigning files to metadata nodes based on their inode number. This distributes metadata operations uniformly and eliminates the metadata bottleneck that affects centralized MDS architectures.

The \textbf{striping architecture} in BeeGFS is fully configurable per directory or file: data is striped across a configurable number of storage targets at a configurable chunk size. For large sequential workloads, a larger chunk size (512KB, 1MB) reduces metadata overhead. For small random I/O, a smaller chunk size may reduce wasted space.

BeeGFS's key differentiator from Lustre and GPFS is its \textbf{ease of deployment and operation}. Installation is straightforward on standard Linux distributions, the management interface is web-based, and the system can be running in a matter of hours on commodity server hardware. This accessibility has made BeeGFS popular in mid-tier HPC installations, enterprise AI/ML clusters, and manufacturing environments that need high-throughput shared storage without dedicated HPC storage administrators.

BeeGFS provides an optional high-availability layer through buddy mirroring: metadata and data targets can be paired, with changes synchronously mirrored between the pair. If one target fails, its buddy continues serving data. This provides resilience without the complexity of external replication management.

\section{Software-Defined Storage: Ceph}
Figure~\ref{fig:ceph-rados} shows the Ceph RADOS data path from client to OSD replicas via placement groups and the CRUSH algorithm.

\begin{figure}[htbp]
\centering
\begin{tikzpicture}[scale=0.85, every node/.style={transform shape}]
  % Client
  \node[component dark] (client) at (0,0) {Client};

  % librados
  \node[component] (librados) at (2.8,0) {librados};

  % hash
  \node[component accent, minimum width=1.8cm] (hash) at (5.4,0) {hash(obj)};

  % PG
  \node[component, minimum width=1.6cm] (pg) at (7.8,0) {PG};

  % CRUSH
  \node[component accent, minimum width=1.8cm] (crush) at (10.2,0) {CRUSH\\map};

  % OSDs
  \node[component green, minimum width=1.6cm] (osd1) at (13.0, 1.0) {OSD.1\\{\tiny primary}};
  \node[component gray,  minimum width=1.6cm] (osd2) at (13.0, 0.0) {OSD.2\\{\tiny replica}};
  \node[component gray,  minimum width=1.6cm] (osd3) at (13.0,-1.0) {OSD.3\\{\tiny replica}};

  % Arrows
  \draw[dataarrow] (client.east) -- (librados.west);
  \draw[dataarrow] (librados.east) -- (hash.west);
  \draw[dataarrow] (hash.east) -- (pg.west);
  \draw[dataarrow] (pg.east) -- (crush.west);
  \draw[dataarrow] (crush.east) -- ++(0.5,0) |- (osd1.west);
  \draw[dataarrow thin] (crush.east) -- ++(0.5,0) |- (osd2.west);
  \draw[dataarrow thin] (crush.east) -- ++(0.5,0) |- (osd3.west);
\end{tikzpicture}
\caption{Ceph RADOS data path: object hashing, placement group mapping, and CRUSH-based OSD selection.}
\label{fig:ceph-rados}
\end{figure}

Ceph is arguably the most architecturally significant distributed storage system of the past fifteen years, representing a complete reimagining of how storage should work at scale. Developed by Sage Weil for his doctoral dissertation at UC Santa Cruz and subsequently commercialized through Inktank (later acquired by Red Hat), Ceph provides block storage, file storage, and object storage from a single unified storage cluster, all built atop a common distributed object store called RADOS.

\section{RADOS: The Reliable Autonomic Distributed Object Store}
RADOS is the foundational layer of Ceph --- a distributed object store that manages its own data placement, replication, and recovery without any centralized controllers. Understanding RADOS is essential to understanding every layer above it.

The RADOS cluster consists of two types of daemons:

\textbf{Object Storage Daemons (OSDs)} are the workhorses of the cluster. Each OSD manages one storage device (typically one OSD per physical drive, though NVMe devices may justify multiple OSDs per drive to saturate their parallelism). OSDs are responsible for storing data objects on their local devices, replicating data to peer OSDs, detecting and reporting failures in peer OSDs through a heartbeat mechanism, and recovering from failures by re-replicating data when a peer OSD goes down. The OSD daemon stack has evolved significantly: the original file-based FileStore backend has been replaced in modern Ceph by \textbf{BlueStore}, which bypasses the host filesystem entirely and manages the block device directly. BlueStore provides checksumming at the object level, inline compression, and significantly improved write performance (particularly for erasure-coded pools) by eliminating the double-write penalty inherent in journaling filesystems.

\textbf{Monitor daemons (MONs)} maintain the authoritative cluster state in the form of a set of cluster maps: the monitor map, OSD map, placement group (PG) map, CRUSH map, and MDS map (for CephFS). Monitors use a Paxos consensus algorithm to maintain quorum and ensure all cluster state changes are agreed upon by a majority before being committed. An odd number of monitors (typically three or five) is required for quorum. The critical insight about monitors is that \textbf{they do not participate in the data path} --- they do not proxy I/O or hold data. Their role is purely coordination and cluster state management.

\section{CRUSH: Controlled Replication Under Scalable Hashing}
CRUSH (Controlled Replication Under Scalable Hashing) is the algorithm by which Ceph computes data placement without lookup tables. It is one of Ceph's most distinctive and important innovations.

When a client wants to store or retrieve a named object, the sequence of operations is:

\begin{itemize}
  \item 1. The object name is hashed to determine its \textbf{Placement Group (PG)} number. PGs are an intermediate abstraction between individual objects (which can number in the billions) and OSDs. A pool is divided into a configurable number of PGs (typically hundreds to thousands), and each object maps to exactly one PG.
2. The PG number is fed into the \textbf{CRUSH algorithm} along with the current \textbf{CRUSH map} to deterministically compute which OSDs should hold the replicas (or erasure code fragments) for that PG.
3. The client contacts the \textbf{primary OSD} for that PG directly, with no intermediary proxy or broker.
\end{itemize}

The CRUSH map is a hierarchical description of the physical topology of the cluster. The hierarchy consists of \textbf{devices} (individual OSDs) at the leaves and \textbf{buckets} at the internal nodes. Default bucket types include host, rack, row, room, datacenter, and zone, though custom types can be defined. Each device has a \textbf{weight} proportional to its storage capacity. CRUSH \textbf{rules} define placement policies in terms of this hierarchy --- for example, "place three replicas, one per distinct rack" --- and the CRUSH algorithm traverses the hierarchy, making pseudo-random selections at each level according to the weights, to produce the final OSD list.

The principal benefit of CRUSH over lookup table-based placement is that the placement computation is \textbf{deterministic and distributed}: any client or OSD can independently compute where any object should be stored, without consulting a central directory. This eliminates the lookup table as a performance bottleneck and single point of failure. The trade-off is that adding or removing OSDs (which changes weights in the CRUSH map) triggers data migration as some PGs are remapped to different OSDs.

To control the extent of data migration during cluster changes, Ceph supports \textbf{weight sets} and the balancer module, which can adjust the effective weights to minimize unnecessary data movement. CRUSH tunable profiles have evolved over Ceph's history --- from argonaut through bobtail, firefly, hammer, to jewel --- progressively improving the quality of data distribution and reducing unnecessary PG remapping when OSDs fail.

\section{Placement Groups: Balancing Granularity and Overhead}
The placement group layer exists because mapping billions of objects directly to OSDs through CRUSH would result in excessive recomputation and tracking overhead, while maintaining one PG per object would mean millions of PGs each with its own replication state machine. PGs represent a sweet spot: each pool has enough PGs to distribute data evenly across OSDs (a common target is 100--200 PGs per OSD), but few enough that the cluster can track and manage their state efficiently.

Every PG has a state (active, clean, degraded, recovering, undersized, etc.) tracked by the monitors, and each PG's primary OSD is responsible for coordinating reads, writes, and recovery for all objects in that PG. The total number of PGs in a cluster must be tuned based on cluster size --- too few PGs results in uneven distribution; too many consumes excessive memory on OSDs and monitors.

\section{CephFS: The Metadata Overlay}
CephFS provides a POSIX-compliant filesystem on top of RADOS. It uses RADOS for data storage, storing file data as RADOS objects, and adds a \textbf{Metadata Server (MDS)} layer for the filesystem namespace. The MDS manages directory trees, inodes, and file metadata, serving filesystem metadata requests from clients.

CephFS's key architectural advancement over simpler parallel filesystems is \textbf{dynamic subtree partitioning}: the MDS cluster can distribute different directory subtrees across multiple active MDS nodes, adapting the distribution in response to load. A heavily accessed directory subtree will be migrated to a less loaded MDS, or split across multiple MDS nodes if load warrants. This provides effective load balancing for metadata workloads without requiring manual partition planning.

CephFS clients use a kernel module or FUSE-based client, and the client caches both data and metadata aggressively. Cache coherence is maintained through \textbf{capabilities}: the MDS grants clients capabilities for specific operations on specific inodes, and must revoke those capabilities before granting conflicting access to other clients. The client's read capability includes a data version, allowing the client to serve reads from cache without consulting the MDS for every operation.

Multiple MDS daemons can be configured, with some active (serving metadata) and others in standby (ready to take over if an active MDS fails). In larger deployments, multiple active MDS daemons operate simultaneously, each responsible for a portion of the namespace.

\section{Ceph Block Device (RBD)}
RADOS Block Device presents block storage on top of RADOS by striping a virtual block device across multiple RADOS objects, typically at 4MB per object. RBD volumes integrate naturally with Linux (via the librbd library and the rbd kernel module), QEMU/KVM, and OpenStack Cinder and Nova. Snapshots are implemented as copy-on-write clones at the RADOS object level --- a snapshot captures the set of RADOS objects as of a point in time, and subsequent writes go to new objects or new object generations. Layering (cloning a snapshot to create a new thin-provisioned RBD image sharing objects with its parent) is fundamental to efficient VM provisioning in cloud environments.

RBD supports both replication pools and erasure-coded data pools (with a replicated metadata pool), the latter providing better storage efficiency for capacity-heavy workloads at the cost of increased CPU overhead and higher write latency.

\section{Ceph Object Gateway (RGW)}
RGW provides S3-compatible and Swift-compatible object storage interfaces on top of RADOS. It functions as a full S3 API implementation: buckets, objects, multipart uploads, versioning, lifecycle policies, server-side encryption, access control lists, and bucket policies are all implemented. RGW translates S3 API operations into RADOS operations, storing object data as RADOS objects and maintaining bucket and object metadata in dedicated RADOS pools.

RGW supports multi-site deployment with zone groups and zones, implementing asynchronous replication between geographic sites for disaster recovery and geographic data locality. Active-active multi-site configurations allow clients at different geographic locations to write to the nearest site, with data subsequently synchronized across all sites.

\section{GlusterFS}
GlusterFS, originally developed by Gluster Inc. and later acquired by Red Hat, takes a different architectural philosophy from Ceph. Rather than building a new distributed object store and layering protocols on top, GlusterFS aggregates existing filesystems --- treating directories on standard Linux filesystems as \textbf{bricks} --- into larger volumes through a translator architecture.

A GlusterFS volume is a logical collection of bricks, assembled through a stack of user-space translators. The translator stack is the core abstraction: each translator is a shared library implementing filesystem operations (create, read, write, stat, etc.) and passing calls up or down the stack. I/O from clients passes through the translator stack on the client side, and is routed to the appropriate brick (or bricks) based on the translator logic.

The two most important cluster translators are:

\textbf{DHT (Distributed Hash Table)}: DHT distributes files across bricks by hashing each file's name to a 32-bit value and mapping that value to a brick or brick subvolume. Each brick is assigned a contiguous range within the 32-bit hash space, covering the full space without gaps or overlaps. A file is stored on exactly the brick whose range contains the file's hash value --- distribution is a routing function, not striping. When a brick is added or removed, hash ranges are redistributed, and files may need to be migrated to their new target bricks through a rebalance operation.

\textbf{AFR (Automatic File Replication)}: AFR implements synchronous replication across a set of bricks, functioning as a distributed RAID-1. Writes are fanned out to all bricks in the replica set before being acknowledged. Reads are served from any brick in the set. AFR uses a change log-based self-healing mechanism: each brick tracks pending changes in extended attributes, and a self-heal daemon runs in the background to detect and repair inconsistencies between replicas.

These translators can be combined: a \textbf{distributed-replicated} volume first replicates files across N bricks (AFR) and then distributes files across replica groups (DHT), providing both redundancy and distribution. A \textbf{dispersed} volume uses erasure coding across bricks, providing storage efficiency comparable to RAID-6.

GlusterFS's greatest strength is its relative simplicity and its ability to reuse existing Linux filesystems (XFS is the strongly preferred underlying filesystem for bricks). Its weaknesses include metadata operations being handled through the DHT translator without a dedicated metadata service --- for operations like directory listings or stat calls, the client must often query multiple bricks --- and suboptimal performance for small-file workloads compared to purpose-built distributed filesystems.

\section{Scale-Out NAS}
Scale-out NAS systems extend traditional NAS capabilities --- POSIX filesystem semantics, NFS/SMB protocol support, rich data management --- to cluster architectures that can grow from tens of terabytes to multiple petabytes within a single namespace. Unlike parallel filesystems designed primarily for HPC bandwidth, scale-out NAS targets the broader enterprise NAS workload: VMware datastores, home directories, content repositories, backup targets, and mixed workloads.

\section{NetApp ONTAP Cluster Mode}
ONTAP's cluster architecture allows up to 24 nodes (for NAS workloads) in a single cluster, all sharing a common namespace. Nodes are organized as HA pairs (two controllers sharing a disk shelf, providing automated failover if one controller fails), and multiple HA pairs are joined into a cluster through a private, dedicated cluster interconnect (historically 10GbE, now 100GbE on high-end systems).

The fundamental virtualization layer in ONTAP is the \textbf{Storage Virtual Machine (SVM)}, formerly called a Vserver. An SVM is a logically separate NAS server with its own namespace, protocol configurations, IP addresses, and authentication context. Multiple SVMs can coexist on the same cluster, each isolated from the others, enabling multi-tenancy. The cluster administrator manages hardware, nodes, and aggregate-level storage, while SVM administrators manage volumes and data within their SVMs.

\textbf{Logical interfaces (LIFs)} are virtual network endpoints associated with SVMs. LIFs are bound to physical network ports on specific nodes but can be migrated to other nodes or ports without interrupting client access --- this is the mechanism for \textbf{LIF migration during failover}: if a node fails, its LIFs migrate to partner node ports automatically. For NAS protocols, this provides transparent failover without client reconnection in most cases.

ONTAP's \textbf{FlexGroup volumes} are the primary scale-out construct: a FlexGroup is a single volume namespace that spans multiple constituent volumes distributed across all nodes in the cluster. Files are placed across constituents using a hash, allowing both namespace capacity and bandwidth to scale with the number of nodes. A FlexGroup can grow to hundreds of petabytes and serve millions of files. ONTAP's \textbf{FlexCache} feature allows a read-cache volume on one cluster to serve data originating on another cluster (or cloud-based ONTAP), reducing latency and bandwidth for geographically distributed workloads.

NetApp's \textbf{ONTAP Mediator} and \textbf{SnapMirror Active Sync} (formerly SnapMirror Business Continuity) provide synchronous replication between clusters for zero-RPO workloads, with the Mediator serving as an external tiebreaker to prevent split-brain during site failures.

\section{Dell PowerScale (Isilon)}
Dell PowerScale, based on the OneFS operating system originally developed by Isilon Systems, implements scale-out NAS through an architecture that combines the filesystem, volume manager, and data protection into a single distributed layer running symmetrically across all nodes in a cluster.

A PowerScale cluster consists of 3 to 252 nodes, each a self-contained server with CPU, memory, networking, and local drives. All nodes run identical OneFS software. \textbf{There is no primary node, no master controller, and no separate metadata server} --- every node is a peer, participating equally in metadata management and data serving. This is the core architectural commitment of OneFS: a single, symmetric distributed filesystem.

The \textbf{OneFS distributed lock manager} coordinates access to files and directories across all nodes, ensuring coherent access when multiple nodes serve requests for the same file simultaneously. File locking uses a token-based mechanism with a global lock coordinator that can migrate in response to load.

Data protection in OneFS is provided by \textbf{FEC (Forward Error Correction)} --- a form of erasure coding native to the filesystem. Rather than RAID applied to physical drives, OneFS stripes file data and parity information across nodes and drives at the file level, using a configurable protection level (N+1, N+2, N+2:1, N+3, etc.). A file protected at N+2 level can tolerate simultaneous failure of any two drives or any two nodes without data loss. This per-file protection level allows mixed protection levels within the same cluster --- critical data can have higher protection while archive data uses more economical settings.

PowerScale nodes are available in several tiers: all-flash (F-series) for performance-intensive workloads, hybrid NVMe-HDD (H-series) for balanced workloads, archive-class (A-series) for cost-effective capacity, and cloud-enabled (B-series with cloud tiering). A single OneFS cluster can mix node types within the node pool hierarchy, tiering data between pools based on access patterns through \textbf{SmartPools} policy-based data movement.

The cluster scales linearly in both capacity and performance: adding nodes adds storage capacity, CPU, memory, and network bandwidth simultaneously. For read-heavy workloads, bandwidth grows nearly linearly with node count. Write performance scales similarly, bounded ultimately by the protection overhead (parity computation).

\section{Qumulo}
Qumulo takes a pragmatic approach to scale-out NAS that prioritizes operational visibility and real-time analytics above all else. Its architecture, based on a shared-nothing cluster of nodes, stores all metadata on flash (SSD) regardless of the underlying media type for data, ensuring consistent metadata performance regardless of cluster load.

The distinguishing feature of Qumulo's design is its \textbf{real-time metadata analytics}: the filesystem is instrumented to continuously aggregate metadata statistics --- file counts, capacity usage, throughput --- bottom-up through the directory tree, without requiring expensive background tree walks. This makes it possible to answer questions like "which directories have grown the most in the past hour?" or "what is the throughput per directory over the past day?" through a web UI or API query, in real time, at petabyte scale. For enterprises dealing with uncontrolled data sprawl, this capability is practically transformative.

Qumulo's \textbf{Scalable Block Store} uses a log-structured approach, organizing data into a protected virtual address space where each 4K block is either a data block or an erasure coding hash block. The ratio of data to parity blocks adjusts dynamically as the cluster grows --- a small cluster uses replication, while a larger cluster transitions to erasure coding for greater efficiency. Erasure coding is performed natively without requiring separate configuration.

Qumulo offers its software in both appliance form (on Qumulo-branded or validated third-party hardware) and as a cloud-native deployment on AWS and Azure, where it uses cloud block storage for data and can scale compute and storage independently. This cloud-native flexibility --- a single Qumulo instance that spans on-premises and cloud nodes --- is a differentiator for hybrid cloud use cases.

\section{VAST Data}
VAST Data represents perhaps the most architecturally radical departure from traditional scale-out NAS design. Founded in 2016, VAST built its platform around three converging technology trends: the falling cost of QLC NAND flash, the availability of NVMe-oF (NVMe over Fabrics) for network-attached storage, and the maturation of storage class memory (SCM/3D XPoint) as a high-endurance write buffer.

VAST's \textbf{DASE (Disaggregated and Shared-Everything) architecture} challenges the assumption that scalable storage must be built from shared-nothing clusters. In VAST's model, compute elements (called \textbf{servers}) are completely stateless --- they run VAST software but hold no persistent data. All persistent data resides in \textbf{JBOF (Just a Bunch of Flash) enclosures} that are NVMe-oF accessible from any server. Any server can access any NVMe device over the fabric.

This disaggregation has several important consequences. First, server failures are non-disruptive --- the server had no authoritative ownership of data; other servers simply pick up the dropped workload. Second, the system can be expanded by adding either servers (for performance) or JBOF capacity (for storage) independently. Third, all servers participate in processing every file, enabling a single large file to be read and processed by all servers simultaneously --- effectively unlimited per-file parallelism.

VAST employs a \textbf{global flash translation layer (GFTL)} that manages all flash devices across all JBOFs as a single pool, applying similarity-based deduplication and compression across the entire namespace. The GFTL organizes data into large erasure-coded stripes of homogeneous data (grouped by similarity hash), dramatically improving compression ratios compared to per-device or per-file compression. This global efficiency engine allows VAST to claim that a single tier of QLC flash can be economically competitive with HDD-based systems for capacity-oriented workloads.

VAST uses \textbf{storage class memory} (3D XPoint/Optane, and more recently, high-endurance flash) as a write buffer across the cluster. Writes land first in the distributed SCM buffer, which provides the high endurance needed to absorb bursty write patterns, and are then organized and flushed to the main QLC flash pool as large, efficient sequential writes. This write path is what allows VAST to use low-cost, low-endurance QLC flash for the bulk of its capacity without wearing out drives prematurely.

VAST's filesystem architecture distributes metadata across all servers, with a globally consistent distributed transaction system ensuring atomicity for multi-object operations. The combination of disaggregated compute, shared flash, and global metadata makes VAST capable of sustaining extremely high concurrent client counts --- thousands of clients accessing billions of files across a single namespace --- which is why VAST has attracted significant adoption in AI/ML training data repositories and large-scale media workflows.

\section{WekaIO (Weka)}
Weka (previously WekaIO) targets the intersection of HPC parallel filesystem performance and enterprise NAS operational simplicity. Its WekaFS is a fully distributed parallel filesystem designed from scratch to operate on NVMe flash, aiming to deliver sub-millisecond latencies across a shared POSIX namespace accessible to thousands of clients simultaneously.

Weka's architecture deploys its storage software inside Linux containers (using cgroups and network namespaces) on commodity server hardware, bypassing the Linux kernel's I/O stack through SPDK (Storage Performance Development Kit). SPDK's user-space NVMe driver eliminates kernel context switches on the I/O path, allowing Weka to approach NVMe device latency directly from the distributed filesystem layer. This bypass architecture is fundamental to achieving the sub-250 microsecond latencies that Weka demonstrates at scale.

Data and metadata are both fully distributed across all nodes. Unlike Lustre's separation of metadata and data servers, Weka treats every node as equivalent --- each node runs both metadata and data services. Metadata is distributed using consistent hashing and protected through replication. Data files are striped across all nodes using a configurable stripe width.

Weka supports all major access protocols from a unified pool: POSIX via its kernel driver or FUSE mount, NFS for compatibility, SMB for Windows clients, S3 for object access, and \textbf{GPUDirect Storage} --- a NVIDIA-developed protocol that allows GPU memory to DMA directly from storage, bypassing the CPU and system memory entirely. GPUDirect Storage support positions Weka as particularly well-suited for AI/ML training workloads where data loading from storage to GPU is often the performance-limiting factor.

Weka can be deployed on-premises on customer hardware, as an SaaS offering on public cloud instances (EC2, Azure VMs, GCE), or in a hybrid configuration where on-premises nodes and cloud instances share a single unified namespace. The cloud integration supports burst-to-cloud for temporary compute scale-out while keeping data on-premises, or full cloud-native deployment for AI training workloads that need high-performance storage adjacent to GPU instances.

\section{Data Placement: Hashing, Consistent Hashing, and CRUSH}
Distributing data across nodes efficiently is a core challenge in any distributed storage system. Several algorithms have been developed, each with different trade-offs in data distribution quality, migration cost, and computational overhead.

\section{Simple Hashing}
The simplest approach is to hash the object key (filename, object name, block address) to a node index using a modular hash: node = hash(key) \% N. This distributes objects uniformly across N nodes and requires only a single computation to locate any object. The fatal weakness is that adding or removing a single node changes N, remapping approximately (N-1)/N fraction of all objects --- effectively causing a full data shuffle across the cluster. For storage systems where data migration is expensive and disruptive, modular hashing is impractical at scale.

\section{Consistent Hashing}
Consistent hashing, introduced by Karger et al. in 1997, addresses the rehashing problem by mapping both nodes and objects to positions on an abstract circular hash ring. Each node is assigned one or more positions on the ring (virtual nodes or vnodes), and each object maps to the nearest node position in the clockwise direction. When a node is added, it takes over responsibility only from its immediate successor on the ring --- approximately 1/N of all objects migrate. When a node is removed, its objects migrate to its successor. The fraction of objects that must move is minimized.

Virtual nodes (multiple ring positions per physical node) smooth the distribution, preventing uneven load when nodes have different capacities or when node counts are small. Both Amazon Dynamo and Apache Cassandra adopted consistent hashing with virtual nodes as their data placement foundation.

For storage systems, consistent hashing provides good distribution and minimal migration but does not natively express failure domain constraints --- there is no built-in mechanism to ensure replicas land on nodes in different racks or datacenters. Layers built on top of consistent hashing must add this logic explicitly.

GlusterFS's DHT translator uses a linear variant of consistent hashing where bricks are assigned ranges on a 32-bit integer space without the circular wrapping, achieving similar distribution properties.

\section{CRUSH: Constraint-Aware Hashing}
CRUSH extends the basic idea of hashing-based placement to incorporate physical topology constraints. Rather than a ring, CRUSH uses a hierarchical tree (the CRUSH map) and a set of placement rules that specify how many replicas to place, at what level of the hierarchy to separate them (rack, row, datacenter), and which device classes to prefer. The CRUSH algorithm traverses the tree from the root, at each level making a weighted pseudo-random selection among the children, guided by a seed derived from the PG number.

The result is placement that is both statistically uniform and constraint-satisfying --- replicas land on different failure domains as specified by the rule. CRUSH is deterministic: given the same CRUSH map and PG number, every node independently computes the same OSD list. No lookup is required.

When the CRUSH map changes (nodes added or removed), CRUSH renames some PG-to-OSD mappings, triggering migration of only the affected data. The balancer module minimizes unnecessary movement by using weight sets that compensate for the probabilistic nature of CRUSH assignments. For very large clusters, even small changes can trigger significant movement if not controlled, which is why operational discipline around CRUSH map changes is important in production Ceph deployments.

\section{Metadata Management in Distributed Systems}
\section{The Metadata Problem}
Filesystem metadata --- directory trees, file attributes, permissions, timestamps, extent maps --- is accessed orders of magnitude more frequently than data in most workloads. A file read involves exactly one data access but may involve multiple metadata operations: path lookup through each directory component, permission check, inode read, extent map lookup. For workloads generating millions of small files (HPC checkpoints, AI training datasets, machine learning feature stores), metadata throughput becomes the dominant constraint.

In a distributed filesystem, metadata management poses a specific challenge: metadata must be consistent across all nodes (a file created on one node must immediately be visible to all other nodes), yet a centralized metadata server becomes a bottleneck and single point of failure as the cluster grows.

\section{Centralized Metadata Architecture}
Lustre's traditional architecture uses a single active MDS (Metadata Server) as the authoritative source for all namespace operations. All clients send metadata requests to the MDS, which serializes them and applies them to the MDT (Metadata Target). The advantages are simplicity --- consistency is trivially maintained since there is one writer --- and predictable behavior for sequential namespace operations.

The disadvantages become apparent at scale. A single MDS can typically sustain 100,000 to 500,000 metadata operations per second depending on workload and hardware. For workloads generating small files at high rates across thousands of compute nodes, this ceiling is quickly reached. Lustre's Distributed Namespace (DNE) partially addresses this by allowing multiple MDTs, with directories assigned to specific MDTs --- but directory striping (introduced in later releases) is required to distribute a single directory's contents across multiple MDTs.

\section{Distributed Metadata Architecture}
Systems like Spectrum Scale (GPFS), BeeGFS, Ceph (with multiple active MDS), and VAST distribute metadata across multiple nodes from the ground up. Each approach has different implementation details:

In \textbf{Spectrum Scale}, metadata is distributed across all nodes that hold NSDs (Network Shared Disks), with file metadata stored in inodes distributed by inode number across the storage pool. A distributed lock manager coordinates access, but no single node is the bottleneck for all metadata operations. The challenge is maintaining consistency across distributed metadata --- a directory create operation must atomically update both the directory's inode and its parent directory's inode, which may be on different nodes.

In \textbf{BeeGFS}, metadata nodes each own a set of directory entries, distributed by directory inode number across all metadata daemons. Each metadata daemon maintains its own local journal, and cross-node operations are handled through a two-phase approach. The uniform distribution ensures that metadata load scales linearly with the number of metadata nodes.

In \textbf{Ceph/CephFS} with multiple active MDS nodes, directory subtrees are dynamically migrated between MDS nodes in response to load. The subtree assignment map is maintained by the MDS cluster itself, with each MDS knowing which subtrees it is authoritative for. Clients discover which MDS to contact for a given path through a request that can be forwarded between MDS nodes if necessary. Dynamic subtree migration allows CephFS to adapt to shifting access patterns without manual reconfiguration.

\section{Metadata Caching and Coherence}
Regardless of whether metadata is centralized or distributed, clients aggressively cache metadata to avoid round-trips to the metadata service on every filesystem operation. Cache coherence --- ensuring that a client's cached metadata remains valid when another client modifies the same file or directory --- is implemented through a variety of mechanisms.

\textbf{Lease-based systems} (NFS v4, CephFS) grant clients time-limited or operation-limited leases on specific metadata. While a client holds a valid lease, it may serve operations from cache without consulting the server. The server can revoke leases to notify clients that their cached data is stale.

\textbf{Inode invalidation} (used in some clustered filesystems) broadcasts invalidation messages to all clients holding a cached copy of an inode when that inode is modified. This is simple but generates O(clients) traffic for popular files.

\textbf{Token/capability systems} (Spectrum Scale, CephFS) grant clients specific capabilities (read, write, create) for specific inodes. Capabilities must be revoked before conflicting access is granted. The metadata server tracks capability grants and can request revocation from a specific client rather than broadcasting to all clients.

The right coherence mechanism depends on the access pattern. For write-mostly workloads with many clients updating different files (the common HPC pattern), aggressive caching with lease-based coherence is effective. For highly shared files accessed by many clients, coherence overhead becomes significant, and applications may need to coordinate access at the application level rather than relying on filesystem coherence.

\section{Security \& Governance}
Distributed and software-defined storage systems introduce a significantly expanded security surface compared to traditional monolithic arrays. Data traverses multiple network paths, resides on potentially many nodes, and is managed by software that must itself be secured and authenticated. A rigorous security posture for distributed storage requires addressing encryption, isolation, and authentication at each layer.

\section{Encryption at Rest in Distributed Systems}
Encrypting data at rest in a distributed filesystem is more complex than in a single-node storage system, because data is spread across many nodes and the encryption key management must be consistent across the cluster.

\textbf{Ceph} supports encryption at the OSD level through the BlueStore backend, which can encrypt all data written to each OSD using dm-crypt (AES-XTS) before it reaches the disk. Keys are typically stored externally in a KMIP-compatible key management server (such as HashiCorp Vault or Thales CipherTrust), integrated through the Ceph key management service (KMS) interface. CephFS can additionally use file-level encryption (fscrypt) at the client level, separating data encryption from storage-level encryption and allowing per-user or per-directory encryption keys.

\textbf{PowerScale/OneFS} supports data-at-rest encryption (DARE) through self-encrypting drives (SEDs) with OneFS managing the key hierarchy using a data authentication key (DAK). Optionally, OneFS integrates with external KMS platforms for key lifecycle management, separating key custody from data custody.

\textbf{Spectrum Scale} offers native filesystem-level encryption (GPFS encryption) that encrypts individual files using AES-256. Encryption policies can be defined per fileset, allowing different encryption domains within the same cluster. Keys are managed through IBM Security Key Lifecycle Manager or Thales, with key derivation hierarchies allowing master keys to protect file encryption keys (FEKs).

\textbf{BeeGFS} does not provide native encryption at rest; it relies on the underlying block device encryption (dm-crypt/LUKS) at each storage target, managed externally.

\textbf{VAST Data} encrypts all data at rest by default, using AES-256-XTS with keys managed through an integrated key management system that can be connected to external KMS through KMIP. Global deduplication across the cluster occurs before encryption (deduplication is performed on plaintext to maintain effectiveness), which has policy implications in multi-tenant environments.

\section{Encryption in Transit}
In distributed storage systems, data moves between clients and storage nodes, between storage nodes for replication, and between the cluster and external systems. Each of these paths requires encryption for deployments in multi-tenant or untrusted-network environments.

\textbf{Ceph} encrypts client-to-cluster traffic using msgr2 (messenger version 2 protocol) with optional CRC or encryption modes. In encryption mode, all traffic between Ceph clients (including OSD-to-OSD replication) is encrypted using AES-GCM. This is enabled per connection type in the Ceph configuration and imposes a CPU overhead on both ends. For high-throughput workloads, hardware AES acceleration (present on all modern server CPUs) is essential to limit this overhead.

\textbf{PowerScale} supports SMB3 encryption for SMB-protocol traffic, NFSv4 with Kerberos for NFS traffic, and TLS for S3 API connections. Internal cluster traffic between nodes is on a dedicated private cluster interconnect network; physical segmentation of this network from external networks is an important architectural control.

\textbf{Spectrum Scale} supports encryption of all inter-node communication using TLS certificates, with per-connection certificate validation. The cluster must be configured with a certificate authority and distributed certificates for each node.

\section{Multi-Tenant Isolation}
Distributed storage systems serving multiple tenants (business units, customers, or application environments) must provide isolation that prevents one tenant's data or operations from affecting another's.

\textbf{Namespace isolation} is the fundamental mechanism. Ceph uses pools as the primary isolation boundary (each pool has independent CRUSH rules, access controls, and encryption policy), with individual users granted access to specific pools through the Ceph authentication system (CephX). CephX uses a shared-secret model (HMAC-SHA1) where each client entity has a set of capabilities specifying which pools and operations it may access.

\textbf{ONTAP SVMs} provide multi-tenant isolation at the NAS layer --- each SVM has its own namespace, network interfaces, and authentication configuration, and an SVM administrator cannot access another SVM's data even on the same cluster. The cluster administrator can set storage quotas, QoS limits, and capacity policies per SVM.

\textbf{PowerScale zones} implement multi-tenant isolation within OneFS, providing separate access zones with independent authentication providers, export configurations, and access policies. Zones share the physical cluster but are logically isolated.

\textbf{Network isolation} through VLAN segmentation or physically separate network interfaces should be combined with namespace isolation for deployments where strong tenant separation is required. Relying solely on access control without network isolation leaves clusters vulnerable to protocol-level attacks that bypass access control.

\section{Certificate-Based Cluster Authentication}
Internal cluster communication --- OSD-to-OSD, OSD-to-monitor, client-to-cluster --- is authenticated using certificate or shared-key mechanisms to prevent rogue nodes from joining the cluster and exfiltrating data or injecting corrupted data.

\textbf{Ceph} uses CephX, a challenge-response authentication protocol based on shared secrets (not certificates). Each daemon and client identity has a key registered with the monitor cluster. When a client connects, the monitors issue time-limited session tickets that encode the client's capabilities. While effective against outsider attacks, CephX provides less protection than certificate-based mutual TLS in environments where key compromise is a concern. Ceph can be deployed behind a TLS-terminating proxy for external access, and msgr2 encryption provides confidentiality for cluster-internal traffic.

\textbf{ONTAP} uses SSL certificates for HTTPS management APIs and cluster peer authentication. iSCSI CHAP and FC-SP provide storage protocol authentication. Cluster-to-cluster peering (for SnapMirror replication) uses per-cluster certificates issued by an internal CA or trusted external CA.

\textbf{PowerScale} requires joining nodes to the cluster through a secure bootstrap process and uses internal certificates for cluster daemon communication. External access is authenticated through Active Directory (for SMB/Kerberos), LDAP, or local user databases.

For all distributed storage systems, hardening recommendations include: limiting cluster management interfaces to dedicated management networks, enabling audit logging for all administrative actions, regularly rotating authentication credentials, integrating with enterprise identity providers (Active Directory, LDAP) rather than local user databases, and conducting periodic access reviews to ensure cluster access follows the principle of least privilege.


\section{Hyperconverged Infrastructure}

Hyperconverged infrastructure emerged as one of the most consequential architectural shifts in enterprise data center design of the 2010s. By collapsing the traditional three-tier model --- separate compute servers, storage arrays, and network interconnects --- into a pool of commodity server nodes in which compute and storage are co-resident and managed as a single system, HCI dramatically simplified deployment, reduced the operational footprint, and brought enterprise storage capabilities to environments previously priced out of dedicated storage arrays.

This chapter examines HCI's architectural foundations, dissects the major HCI platforms in depth, explores the networking prerequisites that make HCI function at enterprise grade, and addresses the scaling models that have emerged as the market matured. Perhaps most importantly, it examines the boundaries of HCI applicability --- the workload characteristics that make HCI an excellent fit, and those where alternative architectures remain superior.

\section{HCI Architecture: The Core Concept}
\section{From Three-Tier to Converged}
The traditional enterprise compute stack is a three-tier architecture: compute servers (often blade or rack servers) connect via Fibre Channel or iSCSI to dedicated storage arrays, while both tiers connect to a converged network fabric. This model optimizes each tier independently --- servers can be scaled for compute, storage arrays for capacity and IOPS, and the network for throughput --- but introduces significant operational complexity. Provisioning a new workload requires coordination across server, storage, and network teams, each with their own toolsets and workflows. Storage arrays require specialized expertise to manage, and capacity planning must account for multiple independent silos.

Converged infrastructure (CI) was the first step toward simplification: vendors like FlexPod (NetApp + Cisco) and Vblock (then VCE, now Dell VxBlock) pre-engineered and pre-validated the combination of servers, storage arrays, and networking into validated reference architectures. Deployment was faster because the integration work was done by the vendor, but the operational model remained three-tier --- storage was still a separate external array.

HCI took the next logical step by eliminating the external storage array entirely. In an HCI system, each node in the cluster contributes its local drives --- SSDs, NVMe devices, and optionally HDDs --- to a software-defined storage pool that is managed by a storage virtualization software layer running on every node. Virtual machine disks and other storage objects are stored on this distributed pool rather than on an external array. A VM running on node A may have its virtual disk data distributed across nodes A, B, and C --- the HCI storage software transparently handles placement, replication, and read/write routing.

\section{The HCI Node}
A standard HCI node is a commodity x86 server that provides CPU, memory, a local storage tier (typically a mix of flash and potentially HDD), and network interfaces. The compute resources run the hypervisor (vSphere, Hyper-V, or KVM, depending on the platform), which in turn runs workload VMs. Alongside these workload VMs, the HCI storage software runs as a privileged VM or service --- in VMware vSAN's case as part of the vSAN kernel module directly in the hypervisor; in Nutanix's case as the Controller VM (CVM) which is a dedicated service VM running on each node.

All storage I/O from workload VMs is intercepted and routed through the HCI storage software, which makes decisions about placement (which nodes hold the data), replication (how many copies or erasure code fragments), and caching (whether to serve from a faster local tier before going to the distributed pool).

The network interconnecting HCI nodes serves a dual purpose: it carries both traditional workload traffic (VM-to-VM, VM-to-external networks) and storage replication traffic (data written on node A being replicated to nodes B and C). This dual-purpose network creates the requirement for adequate bandwidth --- in practice, 10GbE is the minimum for small HCI clusters, with 25GbE or higher strongly preferred for production workloads, and RDMA used for latency-sensitive deployments.

\section{VMware vSAN}
VMware vSAN (Virtual SAN, now commonly referred to as vSAN as part of VMware Cloud Foundation) is the most widely deployed HCI storage platform in enterprise environments, largely because of VMware's dominant position in enterprise virtualization. vSAN is implemented as a kernel module within the ESXi hypervisor, meaning it does not require a separate service VM and has direct access to the I/O path without virtualization overhead.

\section{Architecture Overview}
vSAN clusters consist of a minimum of two ESXi hosts (for a two-node vSAN cluster with a witness) up to 64 hosts per cluster. Each participating host contributes its local storage devices to the shared storage pool. vSAN organizes these devices into a \textbf{disk group} per host: each disk group consists of exactly one cache device (an NVMe drive or high-endurance SSD) and one to seven capacity devices. The cache device absorbs write traffic (acting as a write buffer) and serves hot read data, while the capacity devices provide bulk storage.

In the \textbf{Original Storage Architecture (OSA)}, which preceded vSAN 8.0, write I/O is absorbed into the cache device first, then drained to the capacity devices asynchronously. Read hits on the cache tier are served locally without network traffic. Cache misses trigger a read from the capacity tier, potentially from a remote node.

\section{vSAN Express Storage Architecture (ESA)}
vSAN 8.0 introduced the \textbf{Express Storage Architecture (ESA)}, a major redesign of the I/O stack specifically for all-NVMe configurations. ESA eliminates the discrete cache/capacity device group concept: all NVMe devices participate in a unified pool, and vSAN itself manages a log-structured write path using compression-before-write. By compressing data at admission and writing compressed objects sequentially, ESA improves write efficiency on NVMe drives, reduces write amplification, and enables higher effective RAID 5 and RAID 6 (erasure coding) performance at smaller cluster sizes than OSA could support.

ESA supports RAID 5 with as few as four hosts (compared to OSA's six-host minimum for RAID 5), making erasure coding economically practical for smaller clusters. ESA also improves space efficiency through always-on compression without a performance penalty, as modern NVMe drives sustain sequential writes efficiently. ESA requires all-NVMe storage across all nodes and is not backward compatible with OSA clusters --- it represents a clean break in architecture requiring fresh deployment.

\section{Storage Policy-Based Management (SPBM)}
vSAN's most important operational innovation is \textbf{Storage Policy-Based Management (SPBM)}. Rather than manually configuring RAID levels, replica counts, and placement constraints per volume, administrators define named storage policies that specify these parameters declaratively. When a VM is provisioned, its disks are assigned a storage policy, and vSAN's distributed RAID engine enforces the policy continuously.

Key policy parameters include:

\begin{itemize}
  \item • \textbf{Failures to Tolerate (FTT)}: Specifies how many simultaneous host, disk, or fault domain failures the storage can survive. FTT=1 requires two copies (RAID 1 / mirroring) or three fragments of erasure coding. FTT=2 requires three copies or five fragments.
  \item \textbf{RAID method}: RAID 1 (mirroring) provides better write performance and simpler recovery; RAID 5 or RAID 6 (erasure coding) provides better storage efficiency at the cost of higher write overhead.
  \item \textbf{Checksums}: vSAN can compute and verify end-to-end checksums on all I/O.
  \item \textbf{Encryption}: Per-policy encryption at the vSAN datastore level.
  \item \textbf{IOPS limit}: A maximum IOPS cap per disk or VM.
\end{itemize}

SPBM's power is that policies can be changed at any time without downtime. Changing a VM's storage policy from FTT=1 to FTT=2 triggers vSAN to create an additional replica or recalculate erasure code fragments in the background while the VM continues running. This \textbf{non-disruptive policy enforcement} is a significant operational advantage over traditional storage administration.

\section{Fault Domains}
By default, vSAN treats each host as a failure domain. If a host fails, vSAN considers one copy of any data on that host as lost and begins reprotecting affected components (creating additional copies or recalculating erasure code fragments) on surviving hosts. The Failures to Tolerate setting governs how many such host failures the system can survive.

In real data centers, hosts are often grouped in physical racks with shared power and networking. A rack failure (PDU failure, ToR switch failure) would take down all hosts in the rack simultaneously. \textbf{vSAN Fault Domains} allow administrators to group hosts into named fault domains (representing racks, rows, or rooms), telling vSAN that hosts within the same fault domain should be treated as correlated failure risks. With fault domains defined, vSAN places replicas or erasure code fragments such that each one resides in a different fault domain, protecting against rack-level failures.

\section{Stretched Clusters}
vSAN \textbf{stretched clusters} extend the fault domain concept to geographic distances, distributing a single vSAN cluster across two physical sites with a third tiebreaker site hosting only a witness. The stretched cluster operates as a single vSAN datastore and a single vCenter management domain, with VMs able to run on either site.

Data is written synchronously to both sites --- every write must be acknowledged by both the preferred and non-preferred site before the VM receives an I/O completion. This ensures zero data loss (RPO=0) in the event of a site failure, but requires sufficient inter-site bandwidth and latency constraints. VMware's guidance specifies a maximum round-trip latency of 5ms between data sites for a vSAN stretched cluster (or up to 10ms for some configurations). The witness appliance, deployed at a third site (or in a public cloud), votes in quorum decisions when a site fails to prevent split-brain.

In VCF 9.0, VMware introduced the ability to have disaggregated vSphere compute clusters (separate vSphere clusters consuming storage from a vSAN storage cluster, rather than the traditional HCI topology where compute and storage are on the same hosts) support stretched topologies through site-aware fault domains, extending the stretched cluster model to disaggregated architectures.

\section{vSAN Max: Disaggregated vSAN}
vSAN Max (introduced in vSphere 8.0 Update 1) decouples storage from compute in the vSAN architecture, allowing dedicated storage-only vSAN clusters to serve storage to separate vSphere compute clusters. The storage cluster runs ESXi and vSAN without running workload VMs; the compute cluster runs workload VMs without contributing local storage. Compute hosts access the vSAN storage cluster over the network using NFS or iSCSI.

This disaggregation enables organizations that have adopted HCI to independently scale compute and storage, addressing the constraint of traditional HCI where adding compute capacity necessarily adds storage capacity (often at higher price-per-GB than dedicated storage alternatives). vSAN Max allows a single large storage cluster to serve multiple compute clusters, enabling storage consolidation at the platform level.

\section{Nutanix}
Nutanix was founded in 2009 with the explicit goal of applying hyperscaler web-scale infrastructure design principles to the enterprise data center. The company's core innovation was the Nutanix Controller VM (CVM), a privileged VM running on each node that intercepts all storage I/O from workload VMs and implements the distributed storage layer entirely in software on commodity hardware. Nutanix's AOS (Acropolis Operating System) and the Distributed Storage Fabric built on it represent one of the most sophisticated HCI storage architectures available.

\section{The Controller VM (CVM) Model}
Every Nutanix node runs a dedicated Controller VM alongside workload VMs. The CVM is allocated fixed amounts of vCPU and memory reserved exclusively for storage operations. Workload VM storage I/O is redirected to the local CVM through hypervisor-specific mechanisms (vSphere's VM Storage API for Array Integration, or iSCSI to the CVM for KVM-based deployments). The CVM then handles all storage operations: reading from and writing to local drives, replicating data to remote CVMs over the internal network, serving cached reads, and managing metadata.

The CVM model's principal advantage over the vSAN kernel module approach is \textbf{hypervisor independence}: because the storage layer runs as a VM, it can operate on vSphere, KVM (Nutanix AHV), Hyper-V, and Citrix Hypervisor without requiring hypervisor modifications. Nutanix can evolve its storage stack and deploy updates without coordination with hypervisor vendors. The trade-off is that the CVM consumes resources on every node (typically 12-24 vCPUs and 32-128GB RAM, depending on the deployment) that are not available for workloads.

\section{AOS Distributed Storage Fabric}
The Nutanix storage layer is conceptually organized into several stacking layers, each with distinct responsibilities.

\textbf{OpLog} is the first landing zone for write I/O. When a workload VM writes data, the write is accepted by the local CVM's OpLog --- a write-ahead log maintained on the fastest local storage device (SSD or NVMe). The write is simultaneously replicated to the OpLog of one or two remote CVMs (depending on the configured replication factor --- RF2 or RF3) before being acknowledged to the VM. This synchronous replication at acknowledgment time ensures durability: if the local node fails immediately after acknowledging a write, the remote OpLog copies survive. The OpLog is periodically drained to the Extent Store.

\textbf{Extent Store} is the persistent bulk storage layer, spanning all storage tiers (NVMe, SSD, HDD) across all nodes. Data draining from the OpLog is placed in the Extent Store. Nutanix's \textbf{Intelligent Lifecycle Management (ILM)} monitors access patterns and migrates data between storage tiers based on heat --- frequently accessed extents move to faster tiers, infrequently accessed extents move to slower tiers. The Extent Store uses a log-structured layout on each storage device to maximize sequential write performance.

The fundamental unit of storage in the Extent Store is the \textbf{Extent Group} --- a 1MB or 4MB contiguous block stored as a file on the local storage device. Extents (individual data objects, up to 1MB) are packed into Extent Groups. Multiple extents from different VMs may share an Extent Group, and the Extent Group is the unit of replication and fault tolerance.

\textbf{Curator} is the background process responsible for cluster health, data migration, compression, deduplication, and erasure coding. Curator runs distributed jobs across the cluster to spread load, reorganize data according to tiering policies, and ensure that the configured replication factor is maintained across all data.

\section{Shadow Clones: Distributed Read Caching}
One of Nutanix's most elegant features is the Shadow Clone mechanism, designed to solve the "boot storm" problem in virtual desktop and large server clone deployments.

When many VMs are created from a single base image (as in VDI deployments using linked clones, or in server farms using rapid cloning), they all have a shared read-only base vDisk. Without optimization, all VMs would read from the single node that holds the base vDisk's authoritative copy, creating a hot spot. Shadow Clones detect this multi-reader scenario: when a vDisk is detected to be in "multireader" mode (being read by VMs across multiple nodes), AOS creates read-only shadow copies of the base vDisk on each node that is actively reading it. Subsequent reads from VMs on that node are served from the local shadow copy without generating network traffic.

Shadow Clones are automatically created and managed by AOS. When the base vDisk is modified (ending the multi-reader pattern), shadow copies are invalidated. This mechanism dramatically reduces inter-node network traffic during boot storms and improves VDI read performance to near-local-SSD levels for the base image reads that dominate VDI workloads.

\section{Replication Factor and Fault Tolerance}
Nutanix data is protected by replication factor (RF2 or RF3). RF2 stores two copies of every data block across two different nodes; RF3 stores three copies across three different nodes. The replication factor is set at the container (storage pool partition) level and applies to all VMs in that container.

\textbf{Block Awareness} extends the replication logic to account for multi-node enclosure failures: in a Nutanix block (chassis containing 2-4 nodes), if the block shares power or cooling components, a block-level failure could take all nodes in the chassis offline simultaneously. Block Awareness ensures that data replicas are placed in different blocks (chassis), not just different nodes, so a block-level failure does not result in data loss. Similarly, \textbf{Rack Awareness} ensures replicas span different racks.

Nutanix also supports \textbf{erasure coding} (EC-X) as an alternative to replication for cold data. EC-X is applied post-process to data in the Extent Store that has not been accessed for a configurable period, converting the replicated copies into erasure-coded stripes (typically 4+2 or 8+2). EC-X reduces the storage overhead for cold data (from 100\% overhead with RF2 to \~{}33\% overhead with 4+2 erasure coding) at the cost of more complex recovery and slightly higher read latency for uncached data.

\section{Nutanix Prism and AHV}
Nutanix's management plane is \textbf{Prism}, a unified management interface that presents the cluster as a single system regardless of the number of nodes. Prism Element manages a single cluster; \textbf{Prism Central} provides multi-cluster management, policy governance, and integrated analytics across multiple Nutanix clusters in different sites or cloud environments.

\textbf{AHV (Acropolis Hypervisor)} is Nutanix's own KVM-based hypervisor, offered as an alternative to vSphere, Hyper-V, and Citrix Hypervisor on Nutanix hardware. AHV is included in the Nutanix license without additional cost, which has made it increasingly popular for organizations seeking to reduce or eliminate VMware licensing costs. AHV integrates deeply with AOS storage, enabling features like native VM live migration, AHV-level snapshots coordinated with storage snapshots, and simplified networking through AHV's software-defined networking layer.

\section{Microsoft Azure Stack HCI and Storage Spaces Direct}
Azure Stack HCI is Microsoft's HCI solution, combining commodity server hardware with Windows Server (or the dedicated Azure Stack HCI OS) running Hyper-V for compute virtualization and Storage Spaces Direct (S2D) for software-defined storage. Unlike VMware vSAN and Nutanix, which are primarily on-premises solutions with cloud extensions, Azure Stack HCI is explicitly designed as a hybrid cloud extension of Microsoft Azure --- it is managed through Azure Arc and integrates with Azure services for monitoring, security, and backup.

\section{Storage Spaces Direct Architecture}
Storage Spaces Direct is the software-defined storage layer underlying Azure Stack HCI. It aggregates local drives on all cluster nodes into a unified storage pool managed by the Windows Failover Cluster manager, presenting volumes to Hyper-V VMs and Scale-Out File Server (SoFS) shares.

The S2D architecture consists of several layers:

\textbf{Software Storage Bus}: A virtualized storage fabric implemented in software that makes each node's local drives visible to every other node in the cluster over the cluster interconnect. From the Windows storage stack's perspective, all drives in the cluster appear as if they were locally attached to every node. This abstraction allows the Storage Pool layer to operate without awareness of which physical node holds which drive.

\textbf{Storage Pool}: All drives visible through the Software Storage Bus are aggregated into a single pool (or optionally multiple named pools). The pool maintains a persistent layout of data across all drives and nodes, managing read/write operations and rebalancing capacity.

\textbf{Cache layer}: S2D implements an automatic tiered cache. The fastest drives in each node (NVMe or SAS/SATA SSD) serve as a write buffer and read cache for the slower capacity drives (SATA SSD or HDD). The cache is organized per-node --- each node's fast drives cache that node's local capacity drives. Unlike some HCI platforms that implement a cluster-wide cache, S2D's cache is node-local, which simplifies the cache coherence problem but means that data cached on node A's fast drives cannot directly accelerate reads from node B.

\textbf{Volumes}: Storage Spaces Direct creates volumes (using Resilient File System --- ReFS --- as the underlying filesystem on each volume) with configurable resiliency types:

\begin{itemize}
  \item • \textbf{Two-way mirror}: Two copies across two different nodes. Minimum two nodes; tolerates one node failure.
  \item \textbf{Three-way mirror}: Three copies across three different nodes. Minimum three nodes; tolerates two node failures.
  \item \textbf{Dual parity (erasure coding)}: Reed-Solomon 4+2 or 2+2 erasure coding. Minimum four nodes; provides better storage efficiency than mirroring but requires more nodes.
  \item \textbf{Mirror-accelerated parity}: A hybrid approach where recent writes are stored in a mirrored tier for low write latency, then background-converted to parity encoding for space efficiency.
\end{itemize}

\textbf{ReFS}: All S2D volumes use ReFS (Resilient File System) rather than NTFS. ReFS provides block cloning (copy-on-write at the filesystem level, essential for efficient VM checkpoints and snapshots), integrity streams (checksums on all data), and is optimized for large files. ReFS's block cloning capability is what enables Hyper-V VM checkpoints and clone operations to be nearly instantaneous --- copying a VHDX file on ReFS doesn't copy data; it creates a shared reference until one copy diverges.

\section{Azure Stack HCI and Hybrid Integration}
Azure Stack HCI runs a specialized version of the Windows operating system (the Azure Stack HCI OS) and requires registration with an Azure subscription. After registration, the cluster is managed through Windows Admin Center (on-premises) and Azure Arc (from Azure Portal), giving administrators a single pane of glass regardless of workload location.

Azure Stack HCI integrates with several Azure services natively:

\begin{itemize}
  \item • \textbf{Azure Monitor}: Metrics and logs from all cluster nodes stream to Azure Monitor, providing centralized alerting and dashboard visibility.
  \item \textbf{Microsoft Defender for Cloud}: Security posture assessment and threat detection integrated into the hybrid management experience.
  \item \textbf{Azure Site Recovery}: Replication of Hyper-V VMs from the Azure Stack HCI cluster to Azure for disaster recovery.
  \item \textbf{Azure Backup}: VM-consistent backups to Azure Blob Storage.
  \item \textbf{Azure Arc-enabled data services}: Running Azure SQL Managed Instance or PostgreSQL Hyperscale on the cluster, with Azure-managed operations.
\end{itemize}

The strategic intent is clear: Azure Stack HCI is designed to be the preferred on-premises deployment target for organizations that have standardized on Azure cloud services and want consistent tooling, governance, and security management across on-premises and cloud workloads.

\section{S2D Networking Requirements}
Storage Spaces Direct places stringent requirements on cluster networking because storage replication traffic traverses the same physical network as compute traffic. Microsoft mandates RDMA (Remote Direct Memory Access) for S2D clusters to achieve adequate storage performance and reduce CPU overhead from network processing.

Azure Stack HCI and S2D support two RDMA technologies:

\textbf{RoCEv2 (RDMA over Converged Ethernet v2)}: RDMA over UDP/IP, requiring network switches configured for Data Center Bridging (DCB) with Priority Flow Control (PFC) for lossless transport. PFC is required because RDMA transports cannot handle dropped packets gracefully --- a dropped packet in a standard iSCSI session might be retransmitted, but in an RDMA session, a dropped packet can abort the entire I/O operation. PFC provides per-priority flow control, allowing the network switch to signal the sending NIC to pause transmission on a specific traffic class when buffers are filling, preventing drops.

\textbf{iWARP (Internet Wide Area RDMA Protocol)}: RDMA over TCP/IP. iWARP handles packet loss gracefully (TCP manages retransmission), eliminating the requirement for PFC and DCB. This makes iWARP simpler to deploy in existing Ethernet environments but typically has higher CPU overhead than RoCEv2 for the same bandwidth.

For Azure Stack HCI, a minimum of two network adapters per node dedicated to storage traffic (one for each storage network for redundancy) is recommended, running at 25GbE or higher. Microsoft specifically validates server configurations for Azure Stack HCI (the Azure Stack HCI Catalog) to ensure adequate performance.

\section{HCI Networking: The Fabric That Makes It Work}
HCI networking serves two distinct but interdependent traffic types that must coexist without interfering with each other.

\section{Compute Traffic vs. Storage Traffic}
\textbf{Compute traffic} includes VM-to-VM communication, VM-to-external-network communication, and management traffic. This traffic is latency-sensitive for interactive applications and throughput-sensitive for data-moving workloads, but it generally tolerates modest packet loss (TCP handles retransmission) and variable latency.

\textbf{Storage replication traffic} is fundamentally different in character. Every synchronous write operation generates replication traffic between the local node and one or two remote nodes before the write is acknowledged to the VM. This traffic must be low latency (adding replication latency directly adds to VM I/O latency) and must not be dropped (lost replication traffic either stalls the I/O operation while waiting for retransmission or, in the case of RDMA, aborts the operation outright). Peak storage replication bandwidth can be very high --- in a cluster where every node is performing heavy write workloads, aggregate replication traffic can sustain hundreds of gigabits per second.

For these reasons, 10GbE is widely regarded as too slow for production HCI workloads at scale. \textbf{25GbE has become the baseline for enterprise HCI deployments}, providing sufficient bandwidth for combined compute and storage traffic on dual-NIC configurations with reasonable overhead. 100GbE is increasingly common for large clusters and latency-sensitive workloads.

\section{RDMA in HCI}
RDMA eliminates the CPU from the I/O path for network operations. In a traditional TCP socket transfer, the sender's CPU copies data from the application buffer into the kernel's network buffers, the NIC DMAs data from the kernel buffer onto the wire, and the receiver's CPU handles the interrupt and copies data from the kernel buffer to the application buffer. At 25Gbps or 100Gbps, these CPU copies consume significant core-hours that would otherwise be available for workload VMs.

RDMA allows the NIC to directly read from or write to another server's memory, bypassing both CPUs entirely. The storage replication code constructs an RDMA message specifying the source buffer address and the target server's memory address, and the NIC network stack executes the transfer in hardware. This reduces latency from hundreds of microseconds (for TCP) to tens of microseconds (for RDMA) and frees CPU cores for workload processing.

In HCI contexts, RDMA is used for:

\begin{itemize}
  \item • \textbf{Azure Stack HCI (S2D)}: SMB Direct (SMB3 over RDMA) carries all inter-node storage traffic, including data and metadata replication.
  \item \textbf{vSAN}: vSAN ESA uses RDMA for inter-host data paths in large clusters, reducing storage replication latency.
  \item \textbf{Nutanix}: Supports RDMA for CVM-to-CVM communication in high-performance configurations.
\end{itemize}

RDMA deployment is more operationally complex than standard TCP/IP networking. RoCEv2 requires DCB-capable switches configured correctly with PFC and Enhanced Transmission Selection (ETS) to prioritize storage traffic. Misconfigured PFC can cause priority inversion or deadlock. iWARP is operationally simpler but requires TCP-offload capable NICs (also called RDMA-capable NICs or rNICs) and tends to underperform RoCEv2 in large-scale deployments.

\section{Network Topology Options}
\textbf{Leaf-spine for HCI}: The standard network topology for HCI clusters is a leaf-spine fabric where all HCI nodes connect to leaf switches, and leaf switches connect to spine switches. This provides predictable east-west (node-to-node) latency and bandwidth --- critical for storage replication. Oversubscription ratios between leaf and spine should be kept low (2:1 or 1:1 for storage-heavy environments) to prevent spine congestion.

\textbf{Direct-connect for small clusters}: For small Azure Stack HCI and Nutanix clusters (2-4 nodes), Microsoft and Nutanix support direct-connect topologies where nodes are cross-connected with cables, eliminating switches entirely from the storage network. This provides the lowest possible storage replication latency and simplifies the DCB/PFC configuration, but is only practical for small clusters.

\textbf{Converged vs. dedicated storage networks}: Some organizations run storage and compute traffic on the same physical NICs, using VLAN and QoS policies to prioritize storage traffic (converged topology). Others use dedicated NICs and switches for storage traffic (dedicated topology). Dedicated networks provide better isolation and predictability but require more hardware. The converged approach is common in cost-conscious deployments; dedicated networks are preferred for demanding workloads.

\section{Scaling HCI}
\section{The Coupled Scaling Problem}
Traditional HCI was built on the premise that compute and storage scale together: every time you need more compute capacity, you add a node that also adds storage; every time you need more storage, you add a node that also adds compute. This coupling is elegant in its simplicity --- all nodes are identical, the cluster is homogeneous, and management is uniform.

The problem emerges when workloads have asymmetric resource needs. A data warehouse workload might need 10x more storage than compute, making traditional HCI extremely expensive (you're buying compute you don't need just to add storage). A short-term analytics workload might need 10x the compute but very little additional storage, making it equally wasteful to add full HCI nodes.

\section{Compute-Only Nodes}
All major HCI platforms support \textbf{compute-only nodes}: nodes added to the cluster that contribute CPU and memory but no storage. Workload VMs can run on compute-only nodes and access storage distributed across the storage-contributing nodes. This allows compute to be scaled independently when the compute-to-storage ratio of the workload changes, without purchasing unnecessary storage.

Compute-only nodes are fully supported by vSAN, Nutanix AOS, and Azure Stack HCI. They are a standard part of VDI scale-out strategies, where the base images are stored on storage-contributing nodes and VDI sessions scale out across compute-only nodes.

\section{Storage-Heavy Nodes}
Some workloads --- backup appliances, archive repositories, large media stores --- need more storage than compute. \textbf{Storage-heavy nodes} (also called capacity nodes or high-density storage nodes) carry more drives with higher capacity per node, contributing a disproportionate share of storage relative to compute. Nutanix offers dedicated all-flash and high-capacity storage node SKUs designed for this purpose. vSAN supports non-uniform node configurations within a cluster with appropriate CRUSH-like placement logic.

\section{Disaggregated HCI (dHCI)}
Disaggregated HCI (dHCI) represents an architectural evolution that preserves unified management while separating compute and storage hardware completely. Rather than nodes that run both compute and storage software, dHCI uses dedicated compute servers and dedicated storage servers connected by a high-speed network, with unified management software providing the HCI operational experience.

HPE's dHCI (combining HPE ProLiant servers with HPE Nimble storage), NetApp's HCI offering, and VMware's vSAN Max are examples of this model. The practical advantage of dHCI is that compute and storage can be refreshed on independent hardware lifecycle cycles --- compute nodes can be upgraded when new CPU generations arrive without replacing perfectly good storage infrastructure, and storage can be expanded without adding idle compute capacity.

dHCI's trade-off is slightly higher complexity compared to traditional HCI (there are now two types of infrastructure to manage) and the need for a high-performance, low-latency network between compute and storage nodes. When compute and storage are co-located on the same host, storage I/O traverses the local PCIe bus; when separated to different physical nodes, that I/O traverses the network, which adds latency and consumes network bandwidth.

\section{When HCI Fits and When It Doesn't}
\section{Workloads Where HCI Excels}
\textbf{Virtual Desktop Infrastructure (VDI)} is the canonical HCI workload. VDI is characterized by many small VMs with similar resource profiles, high read-to-write ratios (especially during boot storms, which benefit from Nutanix's Shadow Clones or vSAN caching), and a need for cost-effective, dense deployment. The operational simplicity of HCI --- provisioning a new node to scale out a VDI deployment takes hours, not days --- is compelling for organizations running large VDI environments.

\textbf{General-purpose server virtualization} fits HCI well when workloads are mix of medium-sized VMs running business applications: ERP systems, web servers, development environments, collaboration tools. These workloads are well-served by the storage IOPS and throughput that modern HCI can deliver, and the operational benefits of unified management are significant.

\textbf{Remote and branch office (ROBO)} deployments benefit enormously from HCI because the alternative --- a dedicated storage array and compute servers --- requires on-site storage expertise that most branch offices cannot sustain. A two-node HCI cluster with a cloud-based witness provides reasonable availability with minimal on-site hardware and can be managed centrally.

\textbf{Edge computing} is an extension of the ROBO use case: small, often space-constrained deployments running industrial applications, retail systems, or IoT analytics workloads. HCI vendors have developed compact form factors and ruggedized hardware specifically for this market.

\textbf{Dev/test environments} are well-suited to HCI because they can leverage full VM self-service, snapshot/clone operations are fast and cheap, and the environment can be scaled up or down without architectural changes.

\section{Workloads Where HCI Is a Poor Fit}
\textbf{Large-scale databases with predictable, consistent I/O demand} --- particularly those that can saturate dedicated all-flash arrays --- may not be well served by HCI. The storage replication overhead in HCI consumes resources that a dedicated all-flash array would apply entirely to the workload. For databases requiring more than a few hundred thousand IOPS with sub-millisecond latency, dedicated high-performance storage (such as Pure Storage FlashArray, Dell PowerStore, or NetApp AFF) with dedicated compute typically outperforms HCI.

\textbf{Workloads requiring storage independence}: Large organizations with mature storage teams, sophisticated data management workflows (tiering, thin provisioning at scale, advanced replication topologies), and existing investments in SAN or NAS infrastructure may find HCI's storage feature set limiting. HCI storage is excellent for the VMs running on the platform but does not natively serve external storage consumers (SAN hosts, NAS clients) without additional configuration.

\textbf{Massive-scale storage} with petabytes or tens of petabytes of capacity is generally better served by dedicated scale-out NAS, object storage, or SAN systems. HCI cluster sizes are bounded --- 64 hosts for vSAN, 30 hosts for Nutanix standard clusters --- and the cost-per-usable-terabyte of HCI storage (driven by the compute overhead per node) is higher than purpose-built storage.

\textbf{Applications with asymmetric resource profiles} where compute and storage needs diverge significantly will find HCI's coupled scaling model inefficient. A pure analytics workload that needs large storage but minimal compute, or a pure compute workload (simulation, rendering) that needs negligible storage, will waste resources in a standard HCI deployment.

\textbf{High-performance workloads requiring NVMe-oF}: Applications that can consume NVMe-class storage performance --- genomics pipelines, financial modeling, AI training --- may require dedicated NVMe-oF fabrics or high-performance parallel filesystems that HCI platforms currently cannot match.

\section{Security \& Governance}
HCI systems unify compute and storage in a shared hardware and software environment, which creates specific security requirements. The storage replication network, the hypervisor, the storage management plane, and the data plane all reside in proximity, and compromise of any one layer can potentially affect the others.

\section{Encryption at Rest in HCI}
\textbf{vSAN Encryption} is implemented at the vSAN datastore level, encrypting all data before it is written to any local drive. vSAN datastore encryption uses XTS-AES-256 and integrates with an external Key Management Server (KMS) through the KMIP protocol. vCenter connects to the KMS to retrieve data encryption keys (DEKs), which are stored encrypted (wrapped) by the KMS's key encryption key (KEK). If the KMS becomes unavailable, the cluster can continue operating until the next rekey cycle or node addition, at which point KMS access is required. vSAN also supports \textbf{host-level encryption} in addition to or instead of datastore-level encryption, which encrypts data even before it enters vSAN's distributed write path, providing defense in depth.

vSAN's encryption is not deduplication-friendly: data encrypted with unique IVs cannot be deduplicated at the block level, so organizations must choose between encryption and deduplication efficiency. vSAN ESA encrypts after deduplication and compression are applied, preserving efficiency gains.

\textbf{Nutanix Data at Rest Encryption (DARE)} supports two implementation modes. \textbf{Software encryption} (AES-256-XTS) is implemented in the CVM, encrypting all data before it is written to local drives. \textbf{Self-encrypting drive (SED)} support leverages TCG Opal or FIPS 140-2 compliant drives managed through the Nutanix Key Management Service. Like vSAN, Nutanix integrates with external KMIP-compatible KMS platforms (Thales, Vormetric, IBM Security SKLM) for production deployments. Data encrypted with DARE is still eligible for inline deduplication and compression, as these operations are applied before encryption in the data path.

\textbf{Azure Stack HCI} supports BitLocker encryption for all volumes, with keys stored in Active Directory or Azure Key Vault. BitLocker is a volume-level encryption that encrypts entire S2D volumes using AES-256-XTS. Because BitLocker operates at the volume level (below the ReFS filesystem), deduplication within the Storage Pool is not affected --- deduplication occurs before encryption. Key management through Azure Key Vault provides centralized key lifecycle management and audit logging integrated with Microsoft's cloud management plane.

\section{Micro-Segmentation}
HCI environments running many VMs from different business units, applications, or security tiers benefit from \textbf{micro-segmentation}: applying granular network security policies at the individual VM level, preventing lateral movement between VMs even within the same VLAN or subnet.

\textbf{VMware NSX}, commonly deployed alongside vSAN in VMware Cloud Foundation, implements micro-segmentation through distributed firewall rules enforced at the vNIC level of every VM on every host. A VM's traffic is inspected and filtered regardless of which host it is running on, providing consistent policy enforcement even when VMs migrate between hosts. NSX DFW (Distributed Firewall) rules can be based on VM tags, operating system type, application identity (through integration with identity providers), and security group membership --- providing much richer context than traditional perimeter firewalls that operate solely on IP addresses and ports.

\textbf{Nutanix Flow} (now Flow Network Security) provides equivalent micro-segmentation capability within the Nutanix AHV environment, using policy-based rules applied at the CVM level to inspect and enforce traffic between VMs. Flow's policy model uses application-level groupings --- defining which application tiers (web, application, database) can communicate with which others --- rather than requiring explicit IP-based rules.

For organizations running vSAN or Nutanix without dedicated micro-segmentation products, the minimum security control is strict VLAN segregation by security tier: production VMs, DMZ VMs, and management VMs should be on separate VLANs with firewall enforcement between them, even if micro-segmentation is not applied within each tier.

\section{Management Plane Security}
The HCI management plane --- vCenter for vSAN, Prism for Nutanix, Windows Admin Center and Azure Portal for Azure Stack HCI --- is a high-value target because compromising the management plane provides control over all storage and compute resources in the cluster. Management plane security controls should include:

\begin{itemize}
  \item • \textbf{Multi-factor authentication (MFA)} for all administrative access to management interfaces. vCenter, Prism, and Azure Stack HCI all support integration with SAML-based identity providers (Azure AD, Okta, ADFS) that can enforce MFA.
  \item \textbf{Role-based access control (RBAC)} to limit administrative privileges. vSphere RBAC allows fine-grained roles (read-only, VM administrator, datastore administrator, full administrator) to be assigned to individual users or groups. Nutanix Prism supports equivalent role definitions. The principle of least privilege should be applied rigorously --- storage administrators should not have the ability to modify network configurations, and vice versa.
  \item \textbf{Out-of-band management network isolation}: IPMI/iDRAC interfaces, vCenter, and Prism should be accessible only from dedicated management networks or jump hosts, not from production VLANs.
  \item \textbf{Audit logging}: All administrative actions (VM creates, storage policy changes, host additions, key rotations) should be logged to an external SIEM. vCenter, Nutanix Prism, and Azure Stack HCI all support syslog forwarding for audit events.
  \item \textbf{Immutable cluster configuration backups}: Configuration snapshots of vCenter, Nutanix cluster configuration, and S2D storage pool metadata should be stored in an immutable location (air-gapped backup or WORM storage) to enable recovery from a management plane compromise or ransomware event that destroys configuration data.
\end{itemize}

\section{Storage Network Isolation}
The internal storage replication network --- the network over which HCI nodes exchange replicated data --- should be physically or logically isolated from both workload traffic and external networks. This isolation prevents workload VMs from generating traffic that competes with or disrupts storage replication, and prevents external access to the storage replication protocol.

For vSAN, VMware recommends a dedicated vSAN network (a VMkernel port group with vSAN traffic type enabled) on dedicated physical uplinks. For Azure Stack HCI, dedicated storage adapters are required for S2D traffic. For Nutanix, the CVM management network (used for CVM-to-CVM replication) should be on a separate VLAN from VM data networks.

Storage network traffic should carry no sensitive authentication credentials in cleartext. vSAN's inter-host communication is encrypted when vSAN encryption is enabled. Nutanix CVM-to-CVM traffic can be encrypted using the Nutanix IPsec tunnel feature. S2D traffic uses SMB Direct, and SMB3 encryption can be enabled for S2D inter-node traffic to protect storage replication from network eavesdropping on the cluster interconnect.


\section{DRBD --- Distributed Replicated Block Device}

\textbf{DRBD (Distributed Replicated Block Device)} is a Linux kernel module that implements synchronous block-level replication between two or more nodes over a TCP/IP network. DRBD operates at the block layer, presenting a replicated block device (\texttt{/dev/drbdX}) to the local system that can be formatted with any filesystem. Every write to the DRBD device is simultaneously transmitted to the peer node, which writes it to its own backing device. Only when both nodes have acknowledged the write does the operation complete --- this provides synchronous replication with an RPO of zero.

DRBD supports three replication protocols: \textbf{Protocol A} (asynchronous --- write is acknowledged after local write completes, peer write is in-flight), \textbf{Protocol B} (semi-synchronous --- acknowledged after peer has received data into its TCP buffer), and \textbf{Protocol C} (fully synchronous --- acknowledged only after peer has written to stable storage). Protocol C is the most commonly deployed for production use, as it guarantees that both copies are identical at all times.

DRBD integrates with Linux High Availability frameworks (Pacemaker/Corosync) for automated failover. In a typical HA cluster, DRBD replicates the data volume between two nodes, and Pacemaker manages the active/passive resource: it promotes the DRBD resource to primary on the active node, mounts the filesystem, and starts the application. On node failure, Pacemaker promotes the standby node's DRBD to primary and starts the application there.

\section{Longhorn --- Kubernetes-Native Distributed Storage}

\textbf{Longhorn} is a lightweight, cloud-native distributed block storage system originally developed by Rancher Labs (now part of SUSE) and donated to the Cloud Native Computing Foundation (CNCF) as a sandbox project. Longhorn is designed specifically for Kubernetes environments where each node contributes local storage that Longhorn aggregates into a distributed storage pool.

Longhorn's architecture is based on a microservices model: each volume is managed by a dedicated \textbf{Longhorn Engine} (a user-space iSCSI target based on a custom storage controller), with data replicated synchronously across multiple \textbf{replicas} distributed on different nodes. The replication is at the block level, and Longhorn supports configurable replica counts (typically 2 or 3) for fault tolerance. If a replica or node fails, Longhorn automatically rebuilds the missing replica on a healthy node.

Key features include integrated backup to S3-compatible object storage, volume snapshots (which can be used for backup and disaster recovery), recurring snapshot and backup schedules, and a web-based UI for management. Longhorn uses ReadWriteOnce (RWO) access mode, making it suitable for databases and stateful applications but not for shared filesystem workloads. Its simplicity and tight Kubernetes integration make it popular for edge computing and smaller clusters where a full Ceph deployment would be excessive.


\subsection{NetApp ONTAP Select}

\textbf{ONTAP Select} is a software-defined deployment of NetApp ONTAP that runs as a virtual machine on commodity x86 servers, typically on VMware ESXi or KVM hypervisors. ONTAP Select provides the full ONTAP feature set --- including SnapMirror, SnapVault, FlexClone, deduplication, compression, and multiprotocol access (NFS, SMB, iSCSI, FC) --- without requiring dedicated NetApp hardware. This makes it suitable for remote offices, edge deployments, and lab environments where a full FAS or AFF system is not justified.

ONTAP Select uses local DAS (direct-attached storage) --- either HDDs or SSDs --- on the hypervisor host as its backing storage. It supports both single-node configurations (for non-production use) and multi-node HA clusters (two-node or four-node). The HA mechanism uses ONTAP's standard Storage Failover (SFO) model, with NVRAM data replicated between HA partners over a dedicated internode network.

ONTAP Select integrates seamlessly into the broader NetApp ecosystem: it can be a SnapMirror destination for replicating data from on-premises AFF/FAS arrays to remote sites, and it supports the same REST API, System Manager, and BlueXP management interfaces as hardware ONTAP. This positions it as the edge node in the Data Fabric architecture, extending enterprise storage services to locations where a full hardware deployment is impractical.

